\documentclass[10pt,a4paper]{amsart}
\usepackage[margin=19mm]{geometry}
\usepackage{amsmath,amssymb,mathtools}
\usepackage[scr=rsfs,cal=euler]{mathalfa}
\usepackage{xfrac}
\usepackage[unicode,hidelinks,bookmarks=false]{hyperref}
\newcommand{\ie}{i.e.\ }
\newcommand{\cf}{cf.\ }
\newcommand{\resp}{respectively\ }
\newcommand{\Subsection}{\subsection}
\providecommand{\nobreakdash}{\nobreak\hbox{-}}
\setlength{\emergencystretch}{2em}
\multlinegap0pt
\newcommand{\Z}{\mathbf{Z}}
\newcommand{\Spec}{\operatorname{Spec}}
\usepackage{amssymb}
\usepackage{float}
\usepackage[shortlabels]{enumitem}
\usepackage[normalem]{ulem}
\newtheorem{definition}{Definition}
\newtheorem{theorem}{Theorem}
\newcommand{\C}{\mathbf{C}}
\newcommand{\Gal}{\mathrm{Gal}}
\renewcommand{\P}{\mathbb{P}}
\newcommand{\R}{\mathbf{R}}
\newcommand{\minf}{\mathscr C_\infty}
\title{\Large The Absolute Twistor Line and the Geometry of $\overline{\Spec\Z}$}
\author{Alain Connes, Caterina Consani}
\date{Source: arXiv:2609.00299v1 (CC BY 4.0)}
\begin{document}
\maketitle

\noindent\emph{Source and licence.} \Large The Absolute Twistor Line and the Geometry of $\overline{\Spec\Z}$. Version arXiv:2609.00299v1 (CC BY 4.0), distributed by arXiv under the Creative Commons Attribution 4.0 International licence, http://creativecommons.org/licenses/by/4.0/, as recorded in the arXiv licence metadata for this exact version and shown on the abstract page https://arxiv.org/abs/2609.00299v1. Full licence notice: \texttt{licenses/arxiv-2609.00299-CC-BY-4.0.txt}.\par\smallskip

\noindent\textit{Editorial scope.} Counted unit: the article's theorem thm:twistor\_space\_scheme (source0.tex lines 990--1003). The article's complete original proof of it is delivered with it (source0.tex lines 1005--1034, one \texttt{proof} environment of 30 lines). No mathematics is written, repaired or re-derived: the statement and the proof are the article's own lines, in the article's own notation, and no retained line is substituted or re-flowed.  Retained uncounted as the source-local context the proof needs: lines 597--601; lines 709--715; lines 774--830; lines 929--932.  Nothing was omitted from any retained span, so the package carries no omission record.  No retained line contains a citation, so the wrapper carries no bibliography and no bibliography entry.  No retained line contains a figure environment, an image-inclusion command or drawing code, so no image is delivered and the package declares no asset dependency.  Editorial formatting only, in addition: an AMS \texttt{amsart} preamble that restates the article's own declarations and macros used by the retained lines, namely the environments \texttt{definition}, \texttt{theorem} on separate counters, together with the standard packages those lines need.  The retained environments print in this document as Definition 1, Theorem 1 in order of appearance, whereas the article prints the same items with the numbering of its own full document.  The complete native source of the article as distributed in the arXiv e-print for this exact version is bundled unchanged as \texttt{sources/209015/source0.tex}.
\begin{equation}\label{eq:restriction_maps}
J_+\longmapsto J,
\qquad
J_-\longmapsto\epsilon J.
\end{equation}

\begin{equation}\label{eq:restriction_maps1}
T\longmapsto T^{-1},
\qquad
J_+\longmapsto J,
\qquad
J_-\longmapsto\epsilon J.
\end{equation}

\begin{definition}\label{def:alpha_symmetry}
Let
$
\alpha:\mathcal X\longrightarrow \mathcal X
$
be the involutive automorphism defined on the underlying topological space by
interchanging the two closed points,
\[
\alpha(+)= -,\qquad
\alpha(-)= +,
\]
while fixing the generic point
$
\alpha(\eta)=\eta$.

On the generic affine open subset $\Omega$, the induced automorphism of the
structure sheaf
\[
\alpha^*:\mathcal O(\Omega)\longrightarrow\mathcal O(\Omega)
\]
is determined by
\[
\alpha^*(T)=\epsilon T^{-1},
\qquad
\alpha^*(J)=J^{-1}=\epsilon J.
\]
The first formula realizes the signed inversion of the spatial coordinate,
while the second exchanges the two branches determined by the relation
$J^2=\epsilon$. (Note that
$
J(\epsilon J)=\epsilon J^2=\epsilon^2=1$, 
so that
$
J^{-1}=\epsilon J$.)

The generic automorphism extends uniquely to isomorphisms of the affine
charts
\[
\alpha^*_\pm:
\mathcal O(\Omega_\mp)
\longrightarrow
\mathcal O(\Omega_\pm),
\]
compatible with the restriction maps
\eqref{eq:restriction_maps}. Explicitly,
\[
\alpha_+^*(T^{-1})=\epsilon T,
\qquad
\alpha_+^*(J_-)=J_+,
\]
and
\[
\alpha_-^*(T)=\epsilon T^{-1},
\qquad
\alpha_-^*(J_+)=J_-.
\]
\end{definition}

\begin{equation}\label{twop1}
\mathcal X(\C)\cong
\P^1(\C)_1\sqcup\P^1(\C)_2,
\end{equation}

\begin{theorem}
\label{thm:twistor_space_scheme}
The canonical action of the complex conjugation
$
c\in\Gal(\C/\R)
$
on the $\C$-valued points of the equivariant topos
$\minf$ induces the twistor real structure. More precisely,
on the resulting complex projective line the induced real structure is the
fixed-point-free anti-holomorphic involution:
\[
z\longmapsto-\frac1{\bar z}.
\]
\end{theorem}

\begin{proof}
 Since $\alpha$ interchanges the two copies of $\P^1(\C)$ in   the decomposition \eqref{twop1}, the orbit of the action of the symmetries on the complex points of the unfolded space $\mathcal X(\C)$ are parametrized by their intersection with the first copy $\P^1(\C)_1$. 
The complex conjugation replaces $z\times \{\pm i\}_\pm$ by $\bar z\times \{-\pm i\}_\pm$. The first copy $\P^1(\C)_1$ is obtained by gluing
$
\C\times\{i\}_+$
with
$\C\times\{-i\}_-$, using the canonical map
\[
\C^\times \to \left(\C\times\{i\}_+\right)\times \left(\C\times\{-i\}_-\right), \  \ z\mapsto ((z,i),(z^{-1},-i)).
\]
corresponding to the restriction maps \eqref{eq:restriction_maps1}. We can thus parametrize the points of $\P^1(\C)_1$ by the bijection\footnote{This bijection extends \emph{by continuity} to $\C\cup\{\infty\}$ the parametrization of points belonging to the Zariski dense open set $\Omega(\C)$.}
\[
P:\C\cup \{\infty\}\to \P^1(\C)_1, \  z\mapsto z\times \{ i\}_+, \ \infty  \mapsto 0 \times  \{ -i\}_-.
\]
For $z\in \C^\times$, the point $P(z)$ corresponds to the homomorphism 
\[
\phi: \mathcal O(\Omega)\to H\C, \ \phi(T)=z, \ \phi(J)=i.
\]
For $z\in \C^\times$, the orbit of the point $P(z)$ under the geometric symmetry $\alpha$ is (using Definition \ref{def:alpha_symmetry}) the equivalence class:
\[
[P] = \big\{ P, \, \alpha(P) \big\} = \big\{ (z, i), \, (-z^{-1}, -i) \big\}.
\]
Applying standard complex conjugation $c$ to this orbit yields the conjugate equivalence class:
\[
c([P]) = \big\{ (\bar{z}, -i), \, c(-z^{-1}, -i) \big\} = \big\{ (\bar{z}, -i), \, (-\bar{z}^{-1}, i) \big\}.
\]
To determine the induced action on the quotient space (which we may parameterize by the map $P$), we select the representative of the conjugate orbit that lies on $\P^1(\C)_1$. This representative is exactly $(-\bar{z}^{-1}, i)$. 

Thus, the induced action on the spatial coordinate is $z \mapsto -1/\bar{z}$, and it swaps the points $0$ and $\infty$ (as can be seen by \emph{continuity}).
\end{proof}

\end{document}
